\documentclass[12pt]{article}
\usepackage[margin=1in]{geometry}
\usepackage{amsmath,amssymb,amsthm,mathrsfs}
\usepackage{hyperref}
\usepackage{enumitem}
\usepackage{booktabs}
\usepackage{caption}

\theoremstyle{plain}
\newtheorem{theorem}{Theorem}[section]
\newtheorem{proposition}[theorem]{Proposition}
\newtheorem{lemma}[theorem]{Lemma}
\newtheorem{corollary}[theorem]{Corollary}

\theoremstyle{definition}
\newtheorem{definition}[theorem]{Definition}
\newtheorem{remark}[theorem]{Remark}

\numberwithin{equation}{section}

\newcommand{\R}{\mathbb{R}}
\newcommand{\C}{\mathbb{C}}
\newcommand{\Z}{\mathbb{Z}}
\newcommand{\N}{\mathbb{N}}
\newcommand{\HH}{\mathbb{H}}
\newcommand{\br}{\mathrm{br}}
\newcommand{\vis}{\mathrm{vis}}
\newcommand{\kker}{\mathrm{ker}}
\newcommand{\eff}{\mathrm{eff}}
\newcommand{\obs}{\mathrm{obs}}
\newcommand{\LoN}{\mathrm{LoN}}
\newcommand{\red}{\mathrm{red}}
\newcommand{\hid}{\mathrm{hid}}
\newcommand{\dec}{\mathrm{dec}}
\newcommand{\port}{\mathrm{port}}
\newcommand{\crit}{\mathrm{crit}}
\newcommand{\fU}{\mathfrak{U}}
\newcommand{\fI}{\mathfrak{I}}
\newcommand{\fV}{\mathfrak{V}}
\newcommand{\fB}{\mathfrak{B}}
\newcommand{\fA}{\mathfrak{A}}
\newcommand{\fg}{\mathfrak{g}}
\DeclareMathOperator{\Tr}{Tr}
\DeclareMathOperator{\sech}{sech}
\DeclareMathOperator{\spec}{spec}

\begin{document}

%% ============================================================
%%  TITLE
%% ============================================================
\begin{titlepage}
\centering
\vspace*{1cm}

{\LARGE\bfseries Statistical Mechanics as a Projected\\[6pt]
Determinant Bundle on $X(2)$:\\[6pt]
Boltzmann, KMS, and Gibbs from\\[6pt]
Modular Geometry}

\vspace{1.5cm}

{\large Masayoshi Nakamura$^{1,*}$}

\vspace{0.8cm}

{\normalsize
$^1$\,Kitayono Mental Clinic,
338-0002, 4-20-14 Shimoochiai, Chuo-ku,\\
Saitama City, Saitama, Japan.
\texttt{mjolnir501@gmail.com}\\[4pt]
$^*$\,Corresponding author.
}

\vspace{0.8cm}

{\small
\noindent\textbf{Related deposits:}\\
LoNalogy framework (v87.1): \href{https://doi.org/10.5281/zenodo.19115338}{doi:10.5281/zenodo.19115338}\\
Yang--Mills proof (v90): \href{https://doi.org/10.5281/zenodo.19183838}{doi:10.5281/zenodo.19183838}\\
Particle cosmology (v9) and Hubble reconstruction (v2): companion papers
}

\vspace{1.0cm}

\begin{abstract}
\noindent
We prove that every finite-volume Gibbs many-body system is exactly
representable as the first Fourier--holonomy mode of a determinant
bundle on the level-2 modular curve $X(2)=\Gamma(2)\backslash\HH$.
Concretely, for any quantum system $(\mathcal{H},H)$ with
$e^{-\beta H}$ trace class,
\[
\Tr_{\mathcal{H}}\,e^{-\beta H}
= \int_0^1 e^{-2\pi iu}\,
\det\nolimits_F\!\left(1+e^{2\pi iu}\,e^{-\beta H}\right)du,
\]
where the right-hand side is the $u$-projected LoNalogy determinant
fiber with $\beta=2\pi\tau_2$ and $Q=e^{-2\pi\tau_2}$.

The paper establishes five results, each proved from first principles:
\begin{enumerate}[label=(\roman*)]
\item The Kimura--Th\'evenin principle (Schur complement reduction) is
      algebraically identical to Gaussian hidden-mode elimination, making
      LoNalogy a temperature-free generalisation of the coarse-graining
      core of statistical mechanics.
\item The Boltzmann factor $e^{-\beta E}$ is not a postulate but the
      coordinate expression of the $q$-uniformisation
      $Q=e^{-2\pi\tau_2}$ of $X(2)$.
\item The Matsubara frequency ladder and Bose/Fermi statistics arise
      from the spin structure (half-characteristic sectors) of the
      determinant fiber on $X(2)$.
\item The KMS condition is a theorem of Tomita--Takesaki modular theory
      applied to the $X(2)$-parameterised faithful state, not a
      thermodynamic axiom.
\item The Gibbs variational principle follows as a relative-entropy
      corollary of the KMS state.
\end{enumerate}
The universal projected-determinant theorem (vi) then shows that
\emph{every} finite-volume Gibbs system embeds into this structure.
The single remaining open problem is locality-preserving low-rank
compression of this embedding.
\end{abstract}
\end{titlepage}

\setcounter{page}{1}
\tableofcontents
\newpage


%% ============================================================
\section*{Introduction}
\addcontentsline{toc}{section}{Introduction}
%% ============================================================

Statistical mechanics rests on the Boltzmann--Gibbs postulate: the
equilibrium state of a system with Hamiltonian $H$ at inverse
temperature $\beta$ is described by the density operator
\begin{equation}\label{eq:gibbs}
\rho = \frac{e^{-\beta H}}{Z},\qquad
Z = \Tr\,e^{-\beta H}.
\end{equation}
For over 150 years, this has been treated as a foundational axiom,
justified by ergodic hypotheses, information-theoretic arguments
(Jaynes' maximum entropy), or the equivalence of ensembles.  None of
these derivations explains \emph{why} the weight is exponential in the
energy, as opposed to any other monotone decreasing function.

The LoNalogy framework~\cite{LoNalogyv87,YMmassGap,ParticleCosmo}
provides an answer.  The framework is built on the level-2 modular
curve $X(2)=\Gamma(2)\backslash\HH$ and the para-complex unit
$\jmath^2=+1$.  The uniformising parameter of $X(2)$ is the nome
$Q=e^{2\pi i\tau}$, which on the rectangular slice $\tau=i\tau_2$
becomes
\begin{equation}
Q = e^{-2\pi\tau_2}.
\end{equation}
This is already an exponential.  The determinant sector of LoNalogy
uses $Q$ as the fundamental expansion parameter, and the mode weights
are products of factors $(1-Q^n e^{2\pi iu})$.  When read thermally,
these are precisely the Boltzmann--Gibbs weights.

This paper makes this connection rigorous.  We prove that the entire
logical chain
\begin{equation}
X(2) \;\to\; Q=e^{-2\pi\tau_2} \;\to\; e^{-\beta E}
\;\to\; \text{KMS} \;\to\; \text{Gibbs}
\end{equation}
follows from the modular geometry of $X(2)$ and standard mathematical
theorems (Tomita--Takesaki, Fredholm determinants, relative entropy),
with no thermodynamic postulates.

The paper is organised as follows.  Part~I establishes the Schur
complement as the algebraic core shared by LoNalogy and statistical
mechanics.  Part~II derives the Boltzmann factor from the
$q$-uniformisation of $X(2)$.  Part~III derives the Matsubara ladder
and Bose/Fermi statistics from spin structures.  Part~IV derives the
KMS condition from Tomita--Takesaki theory.  Part~V derives the Gibbs
variational principle.  Part~VI proves the universal projected-determinant
theorem.  Part~VII delineates what remains open.


%% ============================================================
\part{Gaussian Reduction Core}
%% ============================================================

\section{The Kimura--Th\'evenin Principle as Hidden-Mode Elimination}\label{sec:KT}

\subsection{Setup}

\begin{definition}[Full quadratic kernel]
Let $x\in\R^{n_V}$ (visible variables) and $y\in\R^{n_I}$ (hidden/idea
variables).  The full quadratic action is
\begin{equation}
S(x,y) = \frac{1}{2}
\begin{pmatrix}x\\y\end{pmatrix}^{\!\top}
\underbrace{\begin{pmatrix}
L_{\vis} & W\\
W^\top & L_{\fI}
\end{pmatrix}}_{\mathcal{K}_{\mathrm{full}}}
\begin{pmatrix}x\\y\end{pmatrix},
\end{equation}
where $L_{\vis}\in\R^{n_V\times n_V}$, $L_{\fI}\in\R^{n_I\times n_I}$
is positive definite, and $W\in\R^{n_V\times n_I}$ is the bridge
coupling.
\end{definition}

\begin{theorem}[Kimura--Th\'evenin principle = Gaussian elimination]\label{thm:KT}
Integrating out the hidden variable $y$ gives
\begin{equation}\label{eq:gausselim}
\int_{\R^{n_I}} e^{-S(x,y)}\,dy
= \left(\frac{(2\pi)^{n_I}}{\det L_{\fI}}\right)^{1/2}
\exp\!\left[-\frac{1}{2}\,x^\top K_{\vis}^{\red}\,x\right],
\end{equation}
where the visible reduced kernel is the Schur complement
\begin{equation}\label{eq:schur}
\boxed{K_{\vis}^{\red} = L_{\vis} - W\,L_{\fI}^{-1}\,W^\top.}
\end{equation}
\end{theorem}

\begin{proof}
Expand the quadratic form:
\begin{equation}
S(x,y) = \frac{1}{2}x^\top L_{\vis}\,x
+ x^\top W\,y
+ \frac{1}{2}y^\top L_{\fI}\,y.
\end{equation}
Complete the square in $y$.  Define the shifted variable
\begin{equation}
y' := y + L_{\fI}^{-1}W^\top x.
\end{equation}
Then
\begin{align}
S(x,y) &= \frac{1}{2}x^\top L_{\vis}\,x
+ x^\top W(y'-L_{\fI}^{-1}W^\top x)
+ \frac{1}{2}(y'-L_{\fI}^{-1}W^\top x)^\top L_{\fI}
(y'-L_{\fI}^{-1}W^\top x)\\
&= \frac{1}{2}x^\top L_{\vis}\,x
+ x^\top Wy' - x^\top WL_{\fI}^{-1}W^\top x\\
&\quad + \frac{1}{2}y'^\top L_{\fI}\,y'
- y'^\top L_{\fI}L_{\fI}^{-1}W^\top x
+ \frac{1}{2}x^\top WL_{\fI}^{-1}L_{\fI}L_{\fI}^{-1}W^\top x.
\end{align}
The cross terms cancel:
\begin{equation}
x^\top Wy' - y'^\top W^\top x = 0.
\end{equation}
The remaining terms are:
\begin{align}
S(x,y) &= \frac{1}{2}x^\top\!\left(L_{\vis}-WL_{\fI}^{-1}W^\top\right)x
+ \frac{1}{2}y'^\top L_{\fI}\,y'.
\end{align}
The first term depends only on $x$; the second depends only on $y'$.
Since $dy=dy'$ (unit Jacobian for a translation), the Gaussian integral
over $y'$ gives
\begin{equation}
\int_{\R^{n_I}} e^{-\frac{1}{2}y'^\top L_{\fI}\,y'}\,dy'
= \left(\frac{(2\pi)^{n_I}}{\det L_{\fI}}\right)^{1/2}.
\end{equation}
Combining yields~\eqref{eq:gausselim} with $K_{\vis}^{\red}$ as stated.
\end{proof}

\begin{remark}[Temperature-free]
No temperature $\beta$ appears in Theorem~\ref{thm:KT}.  The Schur
complement reduction is a purely algebraic operation on the quadratic
kernel.  Statistical mechanics introduces $\beta$ by setting
$\mathcal{K}_{\mathrm{full}}\mapsto\beta\mathcal{K}_{\mathrm{full}}$;
the Schur complement structure is preserved under this scaling.
Therefore, the Gaussian coarse-graining core of statistical mechanics
is a special case of the Kimura--Th\'evenin principle.
\end{remark}

\subsection{Three-sector extension}

\begin{theorem}[Three-sector Schur complement]\label{thm:threesector}
For the LoNalogy three-sector decomposition
$\fU=\fI\oplus\fV\oplus\fB$ with block kernel
\begin{equation}
\mathbb{K} = \begin{pmatrix}
A & X & 0\\
X^* & M & Y^*\\
0 & Y & C
\end{pmatrix},
\end{equation}
where $A$ and $C$ are invertible, the visible effective kernel is
\begin{equation}
\boxed{M_{\eff} = M - X^*A^{-1}X - Y^*C^{-1}Y.}
\end{equation}
The determinant factorises as
$\det\mathbb{K}=\det A\cdot\det C\cdot\det M_{\eff}$.
\end{theorem}

\begin{proof}
Block Gaussian elimination.  First eliminate the $(1,2)$ and $(2,1)$
blocks using $A^{-1}$: the Schur complement of $A$ in the
$(\fI,\fV)$-block gives $\widetilde{M}=M-X^*A^{-1}X$.  Then eliminate
the $(2,3)$ and $(3,2)$ blocks using $C^{-1}$: the Schur complement of
$C$ in the $(\fV,\fB)$-block gives $M_{\eff}=\widetilde{M}-Y^*C^{-1}Y$.
The determinant identity follows from applying the Schur determinant
formula twice:
$\det\mathbb{K}=\det A\cdot\det\begin{pmatrix}\widetilde{M}&Y^*\\Y&C\end{pmatrix}
=\det A\cdot\det C\cdot\det M_{\eff}$.
\end{proof}

\begin{corollary}[LoNalogy reduced visible law]\label{cor:reducedlaw}
The visible effective action in the LoNalogy framework is defined by
\begin{equation}
e^{-\Gamma_{\vis}^{\red}[\Phi]}
= \int e^{-S_{\mathrm{full}}[\Phi,\phi_{\fI},\phi_{\fB}]}
\,D\phi_{\fI}\,D\phi_{\fB}.
\end{equation}
At the Gaussian level, $\Gamma_{\vis}^{\red}[\Phi]
=\frac{1}{2}\Phi^\top M_{\eff}\,\Phi+\mathrm{const}$.
This is the operational content of the statement: ``observable physics
is the Schur complement of the full theory.''
\end{corollary}

\begin{proof}
Direct application of Theorem~\ref{thm:threesector} to the path integral.
\end{proof}

\begin{proposition}[Statistical mechanics as special case]\label{prop:statmech}
In statistical mechanics, the thermal partition function
\begin{equation}
Z(\beta) = \int e^{-\beta H(q_{\vis},q_{\hid})}\,dq_{\hid}
\end{equation}
is obtained by integrating out the ``bath'' or ``hidden'' degrees of
freedom $q_{\hid}$.  For a quadratic Hamiltonian
$H=\frac{1}{2}(q_{\vis},q_{\hid})^\top\mathcal{K}(q_{\vis},q_{\hid})$,
this is exactly the Schur complement operation with
$\mathcal{K}_{\mathrm{full}}\mapsto\beta\mathcal{K}_{\mathrm{full}}$.
The Schur complement is invariant under this scaling:
\begin{equation}
\beta L_{\vis}-\beta W(\beta L_{\fI})^{-1}\beta W^\top
= \beta\!\left(L_{\vis}-WL_{\fI}^{-1}W^\top\right)
= \beta K_{\vis}^{\red}.
\end{equation}
Therefore, the Gaussian coarse-graining of statistical mechanics is a
$\beta$-scaled special case of the Kimura--Th\'evenin principle.
\end{proposition}

\begin{proof}
Direct substitution.  The $\beta$ factors cancel in $L_{\fI}^{-1}$,
leaving the Schur complement unchanged up to an overall $\beta$ scaling.
\end{proof}

\begin{remark}
The reduced effective action $\Gamma_{\vis}^{\red}$ is more primitive
than the free energy $F=-\beta^{-1}\log Z$.  The former is defined
without reference to temperature; the latter requires $\beta$.
In LoNalogy, $\Gamma_{\vis}^{\red}$ comes first.
\end{remark}


%% ============================================================
\part{Boltzmann Factor from $q$-Uniformisation}
%% ============================================================

\section{The Modular Curve $X(2)$ and the Nome}\label{sec:nome}

\begin{definition}[$X(2)$ and the nome $Q$]
The level-2 modular curve is
\begin{equation}
X(2) = \Gamma(2)\backslash\HH,\qquad
\Gamma(2) = \ker\!\left(\mathrm{SL}_2(\Z)\to\mathrm{SL}_2(\Z/2\Z)\right).
\end{equation}
The uniformising parameter is the nome
\begin{equation}
Q = e^{2\pi i\tau},\qquad \tau\in\HH.
\end{equation}
On the rectangular slice $\tau=i\tau_2$ with $\tau_2>0$:
\begin{equation}\label{eq:nomerect}
Q = e^{-2\pi\tau_2}.
\end{equation}
\end{definition}

\begin{definition}[LoNalogy determinant fiber]\label{def:detfiber}
The LoNalogy determinant on $X(2)$ with holonomy variables
$(u,v)\in\R/\Z\times\R$ is~\cite{LoNalogyv87}
\begin{equation}\label{eq:detfiber}
d(\tau;u,v)
= Q^{-6v(v-1)-1/12}\prod_{n\geq 1}
\left(1-Q^{n-v}e^{2\pi iu}\right)
\left(1-Q^{n+v-1}e^{-2\pi iu}\right).
\end{equation}
The half-characteristic sectors are:
\begin{align}
d_{01} &= Q^{-17/12}\prod_{n\geq1}\left(1-Q^{n-1/2}\right)^2,\label{eq:d01}\\
d_{11} &= Q^{-17/12}\prod_{n\geq1}\left(1+Q^{n-1/2}\right)^2,\label{eq:d11}\\
d_{10} &= 2Q^{1/12}\prod_{n\geq1}\left(1+Q^n\right)^2.\label{eq:d10}
\end{align}
\end{definition}


\section{Boltzmann Weight as $q$-Expansion}\label{sec:boltzmann}

\begin{theorem}[Boltzmann factor from $X(2)$]\label{thm:boltzmann}
The mode weight factors in the LoNalogy determinant~\eqref{eq:detfiber}
are Boltzmann factors with respect to the modular inverse temperature
$\beta_X(\tau):=2\pi\tau_2$.

Specifically, define the LoNalogy thermal scale by
\begin{equation}\label{eq:betaX}
\boxed{\beta_X(\tau) := -\log Q = 2\pi\tau_2.}
\end{equation}
Then each factor in the determinant product has the form
\begin{align}
Q^{n-v}e^{2\pi iu}
&= \exp\!\left[-\beta_X\left((n-v)-\frac{2\pi iu}{\beta_X}\right)\right]
= e^{-\beta_X(E_n-\mu_u)},\\
Q^{n+v-1}e^{-2\pi iu}
&= \exp\!\left[-\beta_X\left((n+v-1)+\frac{2\pi iu}{\beta_X}\right)\right]
= e^{-\beta_X(E_n'+\mu_u)},
\end{align}
where $E_n=n-v$ and $E_n'=n+v-1$ are the mode energies, and
$\mu_u=2\pi iu/\beta_X$ is the chemical potential.
\end{theorem}

\begin{proof}
By~\eqref{eq:nomerect}, $Q=e^{-2\pi\tau_2}=e^{-\beta_X}$.  Therefore
\begin{equation}
Q^{n-v} = e^{-\beta_X(n-v)},\qquad
e^{2\pi iu} = e^{-\beta_X\cdot(-2\pi iu/\beta_X)}.
\end{equation}
The product gives
\begin{equation}
Q^{n-v}e^{2\pi iu}
= e^{-\beta_X(n-v)}e^{-\beta_X(-\mu_u)}
= e^{-\beta_X((n-v)-\mu_u)}.
\end{equation}
This is the standard Boltzmann factor $e^{-\beta(E-\mu)}$ with
$E=n-v$ and $\mu=\mu_u$.  The second factor follows identically
with $E'=n+v-1$ and reversed chemical potential sign.
\end{proof}

\begin{corollary}[Exponentiality is not postulated]
The exponential form of the Boltzmann weight is not a thermodynamic
postulate.  It is the coordinate expression of the uniformising
parameter $Q=e^{-2\pi\tau_2}$ of the modular curve $X(2)$.  Any
other functional form (e.g., power-law) would be inconsistent with
the modular geometry.
\end{corollary}

\begin{proof}
The nome $Q$ is defined as $e^{2\pi i\tau}$; this is the standard
uniformisation of the upper half-plane $\HH$ at the cusp $\tau\to i\infty$.
The exponential structure is dictated by the hyperbolic geometry of $\HH$.
\end{proof}

\begin{remark}[Numerical anchor]
At the LoNalogy strict kernel~\cite{YMmassGap},
$\tau_2^{\kker}=1.02785$, giving
$\beta_X^{\kker}=2\pi\times1.02785=6.4581$.
At the self-dual point $\tau=i$: $\beta_X^{\mathrm{sd}}=2\pi$.
These are not free parameters but outputs of the strict kernel
equation $\nu_{\br}(1-k)^2(1+k)=2$.
\end{remark}


%% ============================================================
\part{Matsubara Ladder and Statistics from Spin Structure}
%% ============================================================

\section{Half-Characteristic Sectors}\label{sec:matsubara}

\begin{theorem}[Matsubara frequencies from spin structure]\label{thm:matsubara}
The half-characteristic sectors~\eqref{eq:d01}--\eqref{eq:d10} of the
LoNalogy determinant fiber encode the Matsubara frequency ladders:
\begin{enumerate}[label=(\roman*)]
\item The sectors $d_{01}$ and $d_{11}$, with mode indices $n-\frac{1}{2}$
($n\geq 1$), correspond to the \textbf{fermionic} Matsubara frequencies
\begin{equation}
\omega_n^{(F)} = (2n-1)\pi T = (2n-1)\frac{\pi}{\beta_X},
\qquad n\geq 1.
\end{equation}
\item The sector $d_{10}$, with mode indices $n$ ($n\geq 1$, integer),
corresponds to the \textbf{bosonic} Matsubara frequencies
\begin{equation}
\omega_n^{(B)} = 2n\pi T = 2n\frac{\pi}{\beta_X},
\qquad n\geq 1.
\end{equation}
\end{enumerate}
\end{theorem}

\begin{proof}
\emph{Step 1: Identify the mode ladder.}
In $d_{01}=Q^{-17/12}\prod_{n\geq1}(1-Q^{n-1/2})^2$, each factor
involves $Q^{n-1/2}=e^{-\beta_X(n-1/2)}$.  The exponent $n-1/2$
runs over the half-integer lattice $\{1/2,3/2,5/2,\ldots\}$.
In $d_{10}=2Q^{1/12}\prod_{n\geq1}(1+Q^n)^2$, each factor involves
$Q^n=e^{-\beta_X n}$.  The exponent $n$ runs over the integer lattice
$\{1,2,3,\ldots\}$.

\emph{Step 2: Identify the statistics.}
The sign in the factor determines the statistics:
$(1-Q^{n-1/2})$ corresponds to Fermi--Dirac (exclusion: the minus sign
enforces the Pauli principle in the Fock-space expansion);
$(1+Q^n)$ corresponds to Bose--Einstein (inclusion: the plus sign allows
multiple occupancy).  The mixed sector $d_{11}$ with $(1+Q^{n-1/2})^2$
represents a half-integer bosonic sector (relevant for supersymmetric
extensions, not used in the minimal framework).

\emph{Step 3: Standard identification.}
In thermal field theory, the Matsubara frequencies for a field at
temperature $T=1/\beta$ are
$\omega_n^{(F)}=(2n+1)\pi T$ (fermionic) and
$\omega_n^{(B)}=2n\pi T$ (bosonic).  The fermionic frequencies arise
from antiperiodic boundary conditions on the thermal circle
$[0,\beta)$, which correspond to half-integer mode indices.
The bosonic frequencies arise from periodic boundary conditions,
corresponding to integer mode indices.

The LoNalogy determinant sectors reproduce this structure exactly:
the spin structure (periodic vs.\ antiperiodic boundary conditions on
the fiber) determines whether the mode index is integer or half-integer,
and hence whether the statistics is bosonic or fermionic.
\end{proof}

\begin{remark}
The Matsubara frequency structure is not imposed by hand in LoNalogy.
It is a consequence of the theta-function factorisation of the
determinant on $X(2)$.  The four half-characteristics
$(a,b)\in\{0,1\}^2$ of the theta functions correspond to the four
possible spin structures on the torus, which in turn correspond to
the four sectors (NS--NS, NS--R, R--NS, R--R) of fermion boundary
conditions.
\end{remark}


%% ============================================================
\part{KMS Condition from Tomita--Takesaki Theory}
%% ============================================================

\section{Modular Automorphism and KMS}\label{sec:kms}

\subsection{Algebraic setup}

\begin{definition}[Modular operator]
Let $\mathfrak{A}$ be a von Neumann algebra acting on a Hilbert space
$\mathcal{H}$, and let $\Omega\in\mathcal{H}$ be a cyclic and
separating vector for $\mathfrak{A}$.  The \emph{Tomita operator} $S$
is defined by
\begin{equation}
S\,A\Omega = A^*\Omega,\qquad A\in\mathfrak{A}.
\end{equation}
Its polar decomposition $S=J\Delta^{1/2}$ defines the \emph{modular
operator} $\Delta$ and the \emph{modular conjugation} $J$.  The
\emph{modular Hamiltonian} is
\begin{equation}
H_{\mathrm{mod}} := -\log\Delta.
\end{equation}
\end{definition}

\begin{theorem}[Tomita--Takesaki]\label{thm:TT}
The modular automorphism group
$\sigma_t(A):=\Delta^{it}A\Delta^{-it}$, $t\in\R$,
satisfies:
\begin{enumerate}[label=(\roman*)]
\item $\sigma_t(\mathfrak{A})=\mathfrak{A}$ for all $t$.
\item The state $\omega(A):=\langle\Omega,A\Omega\rangle$ is a
$\beta=1$ KMS state with respect to $\sigma_t$:
for all $A,B\in\mathfrak{A}$, there exists a function
$F_{A,B}(z)$ analytic on the strip $0<\mathrm{Im}\,z<1$ and
continuous on its closure, such that
\begin{equation}
F_{A,B}(t) = \omega(A\sigma_t(B)),\qquad
F_{A,B}(t+i) = \omega(\sigma_t(B)A).
\end{equation}
\end{enumerate}
\end{theorem}

\begin{proof}
This is the standard Tomita--Takesaki theorem; see~\cite{BratteliRobinson}
for a complete proof.
\end{proof}

\subsection{$X(2)$-parameterised KMS}

\begin{definition}[$X(2)$-parameterised state]
Let $(\mathfrak{A},\mathcal{H},\Omega)$ be a faithful state
arising from the LoNalogy reduced visible sector at modular
parameter $\tau\in\HH$.  Define:
\begin{align}
\beta_X(\tau) &:= 2\pi\tau_2,\\
H_X(\tau) &:= \beta_X(\tau)^{-1}H_{\mathrm{mod}}(\tau)
= \frac{H_{\mathrm{mod}}(\tau)}{2\pi\tau_2},\\
\sigma_t^X(A) &:= e^{itH_X}A\,e^{-itH_X}.
\end{align}
\end{definition}

\begin{theorem}[$X(2)$-KMS condition]\label{thm:kms}
The $X(2)$-parameterised faithful state $\omega_{X,\tau}$ is a
$\beta_X(\tau)$-KMS state with respect to $\sigma_t^X$:
\begin{equation}
\boxed{\omega_{X,\tau}\text{ is }\beta_X(\tau)\text{-KMS w.r.t.\ }\sigma_t^X.}
\end{equation}
\end{theorem}

\begin{proof}
By Theorem~\ref{thm:TT}, the state $\omega$ is $\beta=1$-KMS with
respect to the modular automorphism
$\sigma_t^{\mathrm{mod}}(A)=\Delta^{it}A\Delta^{-it}=e^{-itH_{\mathrm{mod}}}Ae^{itH_{\mathrm{mod}}}$.

Now, $H_X=\beta_X^{-1}H_{\mathrm{mod}}$, so
$\sigma_t^X(A)=e^{itH_X}Ae^{-itH_X}
=e^{it\beta_X^{-1}H_{\mathrm{mod}}}Ae^{-it\beta_X^{-1}H_{\mathrm{mod}}}
=\sigma_{t/\beta_X}^{\mathrm{mod}}(A)$.

The KMS condition with respect to $\sigma_t^{\mathrm{mod}}$ at $\beta=1$
states: for all $A,B$, there exists $F_{A,B}(z)$ analytic on
$\{0<\mathrm{Im}\,z<1\}$ with
$F_{A,B}(t)=\omega(A\sigma_t^{\mathrm{mod}}(B))$ and
$F_{A,B}(t+i)=\omega(\sigma_t^{\mathrm{mod}}(B)A)$.

Define $G_{A,B}(z):=F_{A,B}(z/\beta_X)$.  Then $G$ is analytic on
$\{0<\mathrm{Im}\,z<\beta_X\}$ (since $z/\beta_X$ maps this strip to
$\{0<\mathrm{Im}\,(z/\beta_X)<1\}$), and
\begin{align}
G_{A,B}(t) &= F_{A,B}(t/\beta_X)
= \omega\!\left(A\,\sigma_{t/\beta_X}^{\mathrm{mod}}(B)\right)
= \omega\!\left(A\,\sigma_t^X(B)\right),\\
G_{A,B}(t+i\beta_X) &= F_{A,B}(t/\beta_X+i)
= \omega\!\left(\sigma_{t/\beta_X}^{\mathrm{mod}}(B)\,A\right)
= \omega\!\left(\sigma_t^X(B)\,A\right).
\end{align}
This is precisely the $\beta_X$-KMS condition for $\sigma_t^X$.
\end{proof}

\begin{remark}[Logical order]
In standard statistical mechanics, the KMS condition is typically
\emph{derived from} the Gibbs state $\rho\propto e^{-\beta H}$, or
taken as an equivalent axiom.  In LoNalogy, the logical order is
reversed:
\begin{equation}
X(2)\text{ geometry}\;\to\;\omega_{X,\tau}\;\to\;\Delta_\omega
\;\to\; H_{\mathrm{mod}}\;\to\;\text{KMS}.
\end{equation}
The Gibbs state is not assumed; it will be \emph{derived} in the next
section as a consequence of KMS.
\end{remark}

\subsection{Connection to Hawking temperature}

\begin{proposition}[Hawking temperature from KMS]\label{prop:hawking}
In the black hole realisation of LoNalogy~\cite{ParticleCosmo}, the
horizon dictionary gives
$\Omega_H=\kappa_H/\sqrt{7}$, where $\kappa_H$ is the surface gravity
and $\mathcal{V}_S''(0)=7$ is the Schwarzian floor.  The modular
temperature is $T_*=\sqrt{7}/(2\pi)$.  The physical Hawking temperature
is
\begin{equation}
T_H = \Omega_H\,T_* = \frac{\kappa_H}{\sqrt{7}}\cdot\frac{\sqrt{7}}{2\pi}
= \frac{\kappa_H}{2\pi}.
\end{equation}
This is Hawking's 1975 result, derived here as a special case of the
$X(2)$-KMS structure.
\end{proposition}

\begin{proof}
Direct substitution of the horizon dictionary~\cite{ParticleCosmo}
into the $X(2)$-KMS framework.  The modular Hamiltonian in the
black hole sector is $H_{\mathrm{mod}}=2\pi H_X/\kappa_H$, giving
$\beta_{\mathrm{Hawking}}=2\pi/\kappa_H$ and
$T_H=\kappa_H/(2\pi)$.
\end{proof}


%% ============================================================
\part{Gibbs Variational Principle from KMS}
%% ============================================================

\section{Entropy Maximisation as Relative-Entropy Corollary}\label{sec:gibbs}

\begin{theorem}[Gibbs state from KMS]\label{thm:gibbs}
In the finite-dimensional (or quasi-free) reduced sector, the
$\beta_X(\tau)$-KMS state of Theorem~\ref{thm:kms} is represented
by the density matrix
\begin{equation}\label{eq:gibbsstate}
\rho_{X,\tau} = \frac{e^{-\beta_X(\tau)H_X(\tau)}}{Z_X(\tau)},
\qquad
Z_X(\tau) = \Tr\,e^{-\beta_X(\tau)H_X(\tau)}.
\end{equation}
\end{theorem}

\begin{proof}
In finite dimensions, a $\beta$-KMS state for the automorphism
$\sigma_t(A)=e^{itH}Ae^{-itH}$ is a state $\omega$ satisfying
the KMS condition.  By the standard characterisation theorem
(see~\cite{BratteliRobinson}, Theorem~5.3.10), such a state is
unique and given by $\omega(A)=\Tr(\rho_\beta A)$ with
$\rho_\beta=e^{-\beta H}/\Tr e^{-\beta H}$.  Applying this to
$H=H_X(\tau)$ and $\beta=\beta_X(\tau)$ gives~\eqref{eq:gibbsstate}.
\end{proof}

\begin{theorem}[Gibbs variational principle]\label{thm:gibbsvar}
For any state $\rho$ on $\mathcal{H}$ (i.e., any positive
trace-class operator with $\Tr\rho=1$), the von Neumann entropy
$S(\rho):=-\Tr(\rho\log\rho)$ and the mean energy
$E(\rho):=\Tr(\rho H_X)$ satisfy
\begin{equation}\label{eq:variational}
S(\rho)-\beta_X\,E(\rho) \leq \log Z_X,
\end{equation}
with equality if and only if $\rho=\rho_{X,\tau}$.
\end{theorem}

\begin{proof}
Consider the relative entropy of $\rho$ with respect to
$\rho_{X,\tau}$:
\begin{align}
S(\rho\|\rho_{X,\tau})
&:= \Tr\!\left(\rho\!\left(\log\rho-\log\rho_{X,\tau}\right)\right)\\
&= \Tr\!\left(\rho\log\rho\right)
- \Tr\!\left(\rho\log\rho_{X,\tau}\right).
\end{align}
Now, $\log\rho_{X,\tau}=-\beta_X H_X-\log Z_X$, so
\begin{align}
S(\rho\|\rho_{X,\tau})
&= \Tr(\rho\log\rho)+\beta_X\Tr(\rho H_X)+\log Z_X\\
&= -S(\rho)+\beta_X E(\rho)+\log Z_X.
\end{align}
Since relative entropy is non-negative
($S(\rho\|\sigma)\geq 0$ for all states $\rho,\sigma$, with equality
iff $\rho=\sigma$; this is Klein's inequality):
\begin{equation}
-S(\rho)+\beta_X E(\rho)+\log Z_X \geq 0,
\end{equation}
which rearranges to
\begin{equation}
S(\rho)-\beta_X E(\rho) \leq \log Z_X.
\end{equation}
Equality holds iff $\rho=\rho_{X,\tau}$.

The right-hand side equals $S(\rho_{X,\tau})-\beta_X E(\rho_{X,\tau})$
(by direct computation using $\rho_{X,\tau}=e^{-\beta_X H_X}/Z_X$),
confirming that $\rho_{X,\tau}$ uniquely maximises $S-\beta_X E$.
\end{proof}

\begin{corollary}[Why Gibbs?]
The Gibbs state~\eqref{eq:gibbsstate} is not a thermodynamic postulate.
It is the unique entropy maximiser among all states at fixed
$\beta_X(\tau)$, and $\beta_X(\tau)=2\pi\tau_2$ is the modular inverse
temperature of $X(2)$.  Therefore:
\begin{equation}
\boxed{\text{Gibbs} = \text{unique entropy maximiser of the }X(2)\text{-modular KMS state.}}
\end{equation}
\end{corollary}

\begin{proof}
Combine Theorems~\ref{thm:kms}, \ref{thm:gibbs}, and~\ref{thm:gibbsvar}.
\end{proof}


%% ============================================================
\part{Universal Projected-Determinant Theorem}
%% ============================================================

\section{Statement and Proof}\label{sec:universal}

\begin{definition}[Fermionic Fock space]\label{def:fock}
Let $\mathcal{H}$ be a separable Hilbert space.  The fermionic Fock
space is
\begin{equation}
\mathcal{F}_-(\mathcal{H}) = \bigoplus_{n=0}^{\infty}\Lambda^n\mathcal{H},
\end{equation}
where $\Lambda^0\mathcal{H}=\C$ (vacuum sector) and
$\Lambda^n\mathcal{H}$ is the $n$-th exterior power.  The number
operator $N$ acts on $\Lambda^n\mathcal{H}$ as multiplication by $n$.
The second quantisation of a bounded operator $A$ on $\mathcal{H}$ is
\begin{equation}
\Gamma(A) = \bigoplus_{n=0}^{\infty}\Lambda^n A,
\end{equation}
where $\Lambda^n A$ acts on $\Lambda^n\mathcal{H}$ by
$\Lambda^n A(v_1\wedge\cdots\wedge v_n)=Av_1\wedge\cdots\wedge Av_n$.
\end{definition}

\begin{lemma}[Fock-space trace identity]\label{lem:focktrace}
For any trace-class operator $A$ on $\mathcal{H}$ and any $z\in\C$:
\begin{equation}\label{eq:focktrace}
\Tr_{\mathcal{F}_-(\mathcal{H})}\!\left(z^N\,\Gamma(A)\right)
= \det\nolimits_F(1+zA),
\end{equation}
where $\det_F$ denotes the Fredholm determinant.
\end{lemma}

\begin{proof}
By linearity and the spectral theorem, it suffices to prove this for
diagonal $A$ with eigenvalues $\{a_j\}$.  Then $\Gamma(A)$ acts on
$\Lambda^n\mathcal{H}$ with eigenvalues
$\{a_{j_1}\cdots a_{j_n}:j_1<\cdots<j_n\}$.  Therefore
\begin{align}
\Tr_{\mathcal{F}_-}\!\left(z^N\Gamma(A)\right)
&= \sum_{n=0}^{\infty}z^n\sum_{j_1<\cdots<j_n}a_{j_1}\cdots a_{j_n}\\
&= \sum_{n=0}^{\infty}z^n\,e_n(a_1,a_2,\ldots)\\
&= \prod_j(1+za_j)\\
&= \det\nolimits_F(1+zA),
\end{align}
where $e_n$ denotes the $n$-th elementary symmetric function, and the
product formula for elementary symmetric functions is the standard
identity $\sum_n e_n(a)\,z^n=\prod_j(1+za_j)$.  The product converges
absolutely because $A$ is trace class ($\sum|a_j|<\infty$).
\end{proof}

\begin{lemma}[One-particle projector]\label{lem:projector}
The projector onto the one-particle sector $\Lambda^1\mathcal{H}\cong\mathcal{H}$
in the Fock space is
\begin{equation}\label{eq:projector}
P_1 = \int_0^1 e^{2\pi iu(N-1)}\,du.
\end{equation}
\end{lemma}

\begin{proof}
For any eigenvector of $N$ with eigenvalue $n\in\Z_{\geq 0}$:
\begin{equation}
\int_0^1 e^{2\pi iu(n-1)}\,du
= \begin{cases}1 & \text{if }n=1,\\0&\text{if }n\neq 1,\end{cases}
= \delta_{n,1}.
\end{equation}
This is the standard orthogonality of characters of $\Z$:
$\int_0^1 e^{2\pi imu}\,du=\delta_{m,0}$ applied to $m=n-1$.
Therefore $P_1$ projects onto the $N=1$ eigenspace, which is
$\Lambda^1\mathcal{H}\cong\mathcal{H}$.
\end{proof}

\begin{theorem}[Universal projected-determinant theorem]\label{thm:universal}
Let $(\mathcal{H},H)$ be a quantum system with self-adjoint
Hamiltonian $H$, and suppose $e^{-\beta H}$ is trace class for some
$\beta>0$.  (Every finite-volume lattice many-body system satisfies
this condition.)  Set $\beta=\beta_X(\tau)=2\pi\tau_2$.

Define the LoNalogy determinant fiber with holonomy $u\in[0,1)$:
\begin{equation}\label{eq:detbundle}
Z_{\det}^X(\tau,u;H)
:= \det\nolimits_F\!\left(1+e^{2\pi iu}\,e^{-\beta_X(\tau)H}\right).
\end{equation}

Then the Gibbs partition function is the first Fourier--holonomy mode:
\begin{equation}\label{eq:universalthm}
\boxed{Z_H(\tau)
:= \Tr_{\mathcal{H}}\,e^{-\beta_X(\tau)H}
= \int_0^1 e^{-2\pi iu}\,Z_{\det}^X(\tau,u;H)\,du.}
\end{equation}
\end{theorem}

\begin{proof}
Set $A:=e^{-\beta_X H}$ (trace class by assumption) and $z:=e^{2\pi iu}$.

\emph{Step 1: Express the one-particle trace as a Fock-space trace.}
Since $\Lambda^1\mathcal{H}\cong\mathcal{H}$ canonically, and
$\Lambda^1 A=A$, we have
\begin{equation}
\Tr_{\mathcal{H}}(e^{-\beta_X H})
= \Tr_{\Lambda^1\mathcal{H}}(\Lambda^1 A)
= \Tr_{\mathcal{F}_-(\mathcal{H})}(P_1\,\Gamma(A)).
\end{equation}

\emph{Step 2: Insert the Fourier representation of $P_1$.}
By Lemma~\ref{lem:projector}:
\begin{align}
\Tr_{\mathcal{F}_-}(P_1\,\Gamma(A))
&= \Tr_{\mathcal{F}_-}\!\left(
\left(\int_0^1 e^{2\pi iu(N-1)}\,du\right)\Gamma(A)\right)\\
&= \int_0^1 e^{-2\pi iu}\,
\Tr_{\mathcal{F}_-}\!\left(e^{2\pi iuN}\,\Gamma(A)\right)\,du.
\end{align}
The interchange of trace and integral is justified by the dominated
convergence theorem, since $|\Tr_{\mathcal{F}_-}(z^N\Gamma(A))|
=|\det_F(1+zA)|\leq\det_F(1+|z|\cdot|A|)<\infty$ for $|z|=1$ and
$A$ trace class.

\emph{Step 3: Apply the Fock-space trace identity.}
By Lemma~\ref{lem:focktrace} with $z=e^{2\pi iu}$:
\begin{equation}
\Tr_{\mathcal{F}_-}\!\left(e^{2\pi iuN}\,\Gamma(A)\right)
= \det\nolimits_F\!\left(1+e^{2\pi iu}\,e^{-\beta_X H}\right)
= Z_{\det}^X(\tau,u;H).
\end{equation}

\emph{Step 4: Combine.}
\begin{equation}
\Tr_{\mathcal{H}}\,e^{-\beta_X H}
= \int_0^1 e^{-2\pi iu}\,Z_{\det}^X(\tau,u;H)\,du.
\end{equation}
This is~\eqref{eq:universalthm}.
\end{proof}

\begin{corollary}[Generating functional]\label{cor:generating}
For any self-adjoint observable $O$ on $\mathcal{H}$, define the
source-extended partition function
\begin{equation}
Z_H(J;\tau) := \Tr_{\mathcal{H}}\,e^{-\beta_X(\tau)(H-JO)}.
\end{equation}
Then
\begin{equation}
\boxed{Z_H(J;\tau)
= \int_0^1 e^{-2\pi iu}\,
\det\nolimits_F\!\left(1+e^{2\pi iu}\,e^{-\beta_X(\tau)(H-JO)}\right)du,}
\end{equation}
and all thermal correlation functions are generated by
\begin{equation}
\langle O^n\rangle_\tau
= \frac{1}{Z_H(0;\tau)}
\left.\frac{\partial^n}{\partial J^n}Z_H(J;\tau)\right|_{J=0}.
\end{equation}
Therefore, the entire thermodynamics of any finite-volume system is
encoded in the holonomy-projected LoNalogy determinant bundle.
\end{corollary}

\begin{proof}
Replace $H$ by $H-JO$ in Theorem~\ref{thm:universal}.  The
trace-class condition is preserved for $|J|$ sufficiently small
(by perturbation theory for trace-class operators).  The $J$-derivatives
are standard.
\end{proof}


\section{Why a Single Fiber Does Not Suffice}\label{sec:singlefiber}

\begin{proposition}[Single-fiber insufficiency]\label{prop:singlefiber}
In general,
\begin{equation}
\Tr_{\mathcal{H}}\,e^{-\beta H}
\neq \det\nolimits_F\!\left(1+z\,e^{-\beta H}\right)
\end{equation}
for any fixed $z\in\C$.
\end{proposition}

\begin{proof}
Let $\{E_j\}_{j\in J}$ be the eigenvalues of $H$.  Then
\begin{equation}
\det\nolimits_F(1+ze^{-\beta H})
= \prod_j(1+ze^{-\beta E_j})
= \sum_{n=0}^{\infty}z^n\,e_n\!\left(e^{-\beta E_1},e^{-\beta E_2},\ldots\right),
\end{equation}
where $e_n$ is the $n$-th elementary symmetric function.  In particular:
\begin{align}
z^0\text{ coefficient:}&\quad 1,\\
z^1\text{ coefficient:}&\quad \sum_j e^{-\beta E_j}
= \Tr_{\mathcal{H}}\,e^{-\beta H},\\
z^2\text{ coefficient:}&\quad \sum_{j_1<j_2}e^{-\beta(E_{j_1}+E_{j_2})},\\
&\;\;\vdots
\end{align}
The Gibbs partition function is the $z^1$ coefficient, not the full
determinant.  The higher-order terms correspond to multi-particle
sectors in the Fock space and are generically non-zero.  Therefore,
no single value of $z$ gives $Z_H$ without the holonomy projection.
\end{proof}

\begin{remark}[Interpretation]
The holonomy integral $\int_0^1 e^{-2\pi iu}(\cdots)\,du$ extracts
precisely the $z^1=e^{2\pi iu\cdot 1}$ Fourier mode.  This is a
one-particle projection in the Fock space.  The distinction between
``the system embeds into a determinant fiber'' (false) and ``the
system embeds into a determinant \emph{bundle}'' (true, via
holonomy projection) is the key subtlety.
\end{remark}


\section{Connection to the LoNalogy Determinant Sector}\label{sec:connection}

\begin{proposition}[LoNalogy internal consistency]\label{prop:internal}
The determinant fiber~\eqref{eq:detfiber} in LoNalogy has factors of the
form $(1-Q^{n-v}e^{2\pi iu})$.  For a general spectrum $\{E_j\}$, the
universal theorem replaces these with
\begin{equation}
1+e^{2\pi iu}\,e^{-\beta_X E_j}.
\end{equation}
The sign difference (minus in LoNalogy, plus in the universal theorem)
corresponds to the choice of statistics: the LoNalogy determinant uses
fermionic conventions (alternating signs in the exterior algebra), while
Theorem~\ref{thm:universal} uses the standard Fredholm determinant
$\det_F(1+zA)$.  The two are related by $z\mapsto -z$:
\begin{equation}
\prod_j(1-ze^{-\beta E_j}) = \det\nolimits_F(1-zA),
\end{equation}
and the holonomy projection adapts accordingly
($e^{-2\pi iu}\mapsto(-1)e^{-2\pi iu}$, i.e., a half-period shift in $u$).

The roles of $u$ (holonomy / chemical potential) and
$Q=e^{-\beta_X}$ (thermal weight) are the same in both the LoNalogy
determinant sector and the universal theorem.  Therefore,
Theorem~\ref{thm:universal} is not a foreign embedding but an
\emph{internal statement} of the LoNalogy framework.
\end{proposition}

\begin{proof}
Direct comparison of~\eqref{eq:detfiber} and~\eqref{eq:detbundle}.
The holonomy variable $u$ and the nome $Q$ play identical roles.
The sign convention is a matter of the choice of Fock-space grading
(fermionic vs.\ bosonic second quantisation).
\end{proof}


%% ============================================================
\part{No-Go Theorem for Low-Rank Single-Fiber Compression}
%% ============================================================

\section{Statement and Proof}\label{sec:nogo}

The universal projected-determinant theorem (Theorem~\ref{thm:universal})
represents any Gibbs partition function as a holonomy-projected
determinant bundle.  A natural question is whether this can be
improved: can one find a \emph{single} determinant fiber of
lower rank that reproduces the partition function exactly?

The answer is no.

\begin{theorem}[No-go for universal low-rank compression]\label{thm:nogo}
Let $(\mathcal{H},H)$ be a quantum system with $\dim\mathcal{H}=D$
and generic nondegenerate spectrum $\{E_1<E_2<\cdots<E_D\}$.
Suppose there exists an auxiliary Hilbert space $\mathcal{K}$ with
$\dim\mathcal{K}=r$ and a self-adjoint operator $K$ on $\mathcal{K}$
such that
\begin{equation}\label{eq:lowrank}
\Tr_{\mathcal{H}}\,e^{-\beta H}
= \int_0^1 e^{-2\pi iu}\,
\det\nolimits_F\!\left(1+e^{2\pi iu}\,e^{-\beta K}\right)du
\end{equation}
for all $\beta>0$.  Then $r\geq D$.

For generic nondegenerate $H$, equality $r=D$ is required, and the
eigenvalues of $K$ must coincide with those of $H$ (as multisets).
\end{theorem}

\begin{proof}
Let $\{E_a\}_{a=1}^D$ be the eigenvalues of $H$ and
$\{\varepsilon_j\}_{j=1}^r$ be the eigenvalues of $K$.

\emph{Step 1: Extract the first holonomy coefficient.}
By the same argument as in Theorem~\ref{thm:universal}, the
holonomy integral on the right-hand side of~\eqref{eq:lowrank}
extracts the $z^1$ coefficient of $\det_F(1+z\,e^{-\beta K})$:
\begin{equation}
\int_0^1 e^{-2\pi iu}\,\det_F\!\left(1+e^{2\pi iu}\,e^{-\beta K}\right)du
= \sum_{j=1}^r e^{-\beta\varepsilon_j}
= \Tr_{\mathcal{K}}\,e^{-\beta K}.
\end{equation}

\emph{Step 2: Equate exponential sums.}
The hypothesis~\eqref{eq:lowrank} therefore reduces to
\begin{equation}\label{eq:expsum}
\sum_{a=1}^D e^{-\beta E_a} = \sum_{j=1}^r e^{-\beta\varepsilon_j}
\qquad\text{for all }\beta>0.
\end{equation}

\emph{Step 3: Uniqueness of exponential sums.}
Both sides of~\eqref{eq:expsum} are Laplace transforms of discrete
measures:
\begin{equation}
\mu_H = \sum_{a=1}^D \delta_{E_a},\qquad
\mu_K = \sum_{j=1}^r \delta_{\varepsilon_j}.
\end{equation}
The Laplace transform $\mathcal{L}[\mu](\beta)=\int e^{-\beta x}\,d\mu(x)$
is injective on the space of finite Borel measures on $[0,\infty)$
(this follows from the injectivity of the bilateral Laplace transform,
or equivalently from the uniqueness theorem for exponential
polynomials~\cite{Simon}).

Therefore $\mu_H=\mu_K$ as measures, which means the multisets
$\{E_a\}_{a=1}^D$ and $\{\varepsilon_j\}_{j=1}^r$ coincide
(with multiplicities).

\emph{Step 4: Rank bound.}
For generic nondegenerate $H$, the multiset $\{E_a\}_{a=1}^D$ consists
of $D$ distinct elements.  To match this with $\{\varepsilon_j\}_{j=1}^r$,
we need $r\geq D$ (since a set of $r$ elements can contain at most $r$
distinct values).  For exact matching, $r=D$ and $\varepsilon_j=E_{\pi(j)}$
for some permutation $\pi$.
\end{proof}

\begin{remark}[Interpretation]
This theorem shows that the projected-determinant representation
(Theorem~\ref{thm:universal}) is \emph{optimal} in the following sense:
the holonomy projection cannot be eliminated without using an auxiliary
space of full dimension $D$.  The bundle structure (parameterised by the
holonomy $u$) is not a weakness of the representation but an essential
feature.

The distinction is:
\begin{equation}
\boxed{\text{Single fiber: requires }r\geq D.\qquad
\text{Projected bundle: works for any }D\text{ with }r=D.}
\end{equation}
The bundle adds one continuous parameter $u\in[0,1)$ and removes the
need for any rank reduction.
\end{remark}


%% ============================================================
\part{Transfer Class and Locality-Preserving Reduction}
%% ============================================================

\section{The Transfer Class}\label{sec:transfer}

The projected-determinant theorem treats $\mathcal{H}$ as a monolithic
space, ignoring any tensor-product structure.  For lattice many-body
systems, we want \emph{locality-preserving} representations.  The
key concept is the transfer class.

\begin{definition}[Transfer class]\label{def:transfer}
Let $\Lambda=[1,L]\times B$ be a lattice with $L$ sites in the
propagation direction and $B$ a finite transverse section.  A
partition function $Z_\Lambda(\beta)$ belongs to the \emph{exact
transfer class} if there exists a finite-dimensional auxiliary
(boundary) space $\mathcal{B}_\partial$ and a linear operator
$T_\beta:\mathcal{B}_\partial\to\mathcal{B}_\partial$ such that
\begin{equation}\label{eq:transfer}
Z_\Lambda(\beta) = \Tr_{\mathcal{B}_\partial}\,T_\beta^{\,L}.
\end{equation}
The \emph{bridge rank} is $\chi:=\dim\mathcal{B}_\partial$.
\end{definition}

\begin{proposition}[Transfer class = locality-preserving LoN determinant class]\label{prop:transferLoN}
If $Z_\Lambda(\beta)$ belongs to the exact transfer class with
transfer matrix $T_\beta$ and bridge rank $\chi$, then
\begin{equation}
\boxed{Z_\Lambda(\beta)
= \int_0^1 e^{-2\pi iu}\,
\det\nolimits_F\!\left(1+e^{2\pi iu}\,T_\beta^{\,L}\right)du.}
\end{equation}
This is a locality-preserving LoN determinant representation: the
operator $T_\beta$ acts on the boundary space $\mathcal{B}_\partial$
of dimension $\chi$, which scales with the \emph{boundary} of the
system, not its volume.
\end{proposition}

\begin{proof}
Apply Theorem~\ref{thm:universal} with $\mathcal{H}=\mathcal{B}_\partial$
and $H$ replaced by the operator $-\frac{1}{\beta}\log T_\beta$
(or more precisely, note that $\Tr T_\beta^L$ is already of the form
$\Tr A$ with $A=T_\beta^L$, which is the $z^1$ holonomy coefficient
of $\det_F(1+zA)$ by the same Fock-space argument).
\end{proof}


\section{Finite-Range Classical Systems}\label{sec:classical}

\begin{theorem}[Classical local systems are in the transfer class]\label{thm:classical}
Let $(\Lambda,H)$ be a $d$-dimensional finite lattice spin system with
local state space of size $q$ at each site, finite interaction range $R$
in the propagation direction, and transverse section $B$ with $|B|$ sites.

Then there exists an exact transfer matrix $T_\beta$ on a boundary space
$\mathcal{B}_\partial$ with
\begin{equation}\label{eq:chibound}
\dim\mathcal{B}_\partial \leq q^{R|B|}
\end{equation}
such that
\begin{equation}
Z_\Lambda(\beta) = \Tr_{\mathcal{B}_\partial}\,T_\beta^{\,L-R+1}.
\end{equation}
Consequently,
\begin{equation}
\boxed{Z_\Lambda(\beta) = \int_0^1 e^{-2\pi iu}\,
\det\nolimits_F\!\left(1+e^{2\pi iu}\,T_\beta^{\,L-R+1}\right)du.}
\end{equation}
\end{theorem}

\begin{proof}
\emph{Step 1: Decompose the Hamiltonian.}
Write the lattice sites as $(t,\mathbf{x})$ where $t\in[1,L]$ is the
propagation coordinate and $\mathbf{x}\in B$ is the transverse
coordinate.  The Hamiltonian decomposes as
\begin{equation}
H = \sum_{t=1}^{L-R+1} h_t(\sigma_t,\sigma_{t+1},\ldots,\sigma_{t+R-1}),
\end{equation}
where $\sigma_t=(\sigma_{t,\mathbf{x}})_{\mathbf{x}\in B}\in\{1,\ldots,q\}^B$
is the configuration of the $t$-th slice, and $h_t$ depends on at most
$R$ consecutive slices due to the finite interaction range.

\emph{Step 2: Define the boundary state space.}
The auxiliary space is
\begin{equation}
\mathcal{B}_\partial = \C^{q^{R|B|}},
\end{equation}
with basis vectors labelled by the configurations of $R$ consecutive
transverse slices: $|\sigma_t,\sigma_{t+1},\ldots,\sigma_{t+R-1}\rangle$.

\emph{Step 3: Define the transfer matrix.}
The transfer matrix $T_\beta$ acts by
\begin{equation}
\langle\sigma_{t+1},\ldots,\sigma_{t+R}|T_\beta|\sigma_t,\ldots,\sigma_{t+R-1}\rangle
= e^{-\beta h_t(\sigma_t,\ldots,\sigma_{t+R-1})}\cdot
\delta_{\sigma_{t+1}\cdots\sigma_{t+R-1}},
\end{equation}
where the $\delta$ enforces that the overlap region
$(\sigma_{t+1},\ldots,\sigma_{t+R-1})$ is shared between the two
boundary states.  (For $R=1$, this reduces to the standard row-to-row
transfer matrix.)

\emph{Step 4: Verify the trace identity.}
The partition function with periodic boundary conditions is
\begin{align}
Z_\Lambda(\beta) &= \sum_{\{\sigma\}} e^{-\beta H(\sigma)}
= \sum_{\{\sigma\}} \prod_{t=1}^{L-R+1} e^{-\beta h_t}\\
&= \sum_{\sigma_1,\ldots,\sigma_{R}}
\langle\sigma_1,\ldots,\sigma_R|
T_\beta^{L-R+1}
|\sigma_1,\ldots,\sigma_R\rangle\\
&= \Tr_{\mathcal{B}_\partial}\,T_\beta^{L-R+1}.
\end{align}
The LoN determinant representation follows from
Proposition~\ref{prop:transferLoN}.

\emph{Step 5: Reduction is genuine.}
The original state sum has $q^{L|B|}$ terms (volume scaling).
The transfer matrix has dimension $q^{R|B|}$ (boundary scaling).
Since $R$ is fixed and $L\to\infty$, the reduction from volume to
boundary is genuine and not a mere rewriting.
\end{proof}


\section{Example: One-Dimensional Ising Model}\label{sec:ising}

\begin{corollary}[1D Ising in LoN determinant class]\label{cor:ising}
The one-dimensional Ising model with Hamiltonian
\begin{equation}
H(\sigma) = -J\sum_{i=1}^N\sigma_i\sigma_{i+1}
-h\sum_{i=1}^N\sigma_i,\qquad \sigma_i=\pm1,
\end{equation}
(periodic boundary: $\sigma_{N+1}=\sigma_1$) has transfer matrix
\begin{equation}
T = \begin{pmatrix}
e^{\beta J+\beta h} & e^{-\beta J}\\
e^{-\beta J} & e^{\beta J-\beta h}
\end{pmatrix}
\end{equation}
with bridge rank $\chi=2$, and the LoN determinant representation is
\begin{equation}\label{eq:isingLoN}
\boxed{Z_N(\beta)
= \int_0^1 e^{-2\pi iu}\,
\det\!\left(1+e^{2\pi iu}\,T^N\right)du.}
\end{equation}
\end{corollary}

\begin{proof}
\emph{Step 1: Transfer matrix.}
The Hamiltonian decomposes as $H=\sum_{i=1}^N h_i$ with
$h_i=-J\sigma_i\sigma_{i+1}-\frac{h}{2}(\sigma_i+\sigma_{i+1})$.
The transfer matrix element is
\begin{equation}
T_{\sigma_i,\sigma_{i+1}}
= e^{-\beta h_i}
= e^{\beta J\sigma_i\sigma_{i+1}+\frac{\beta h}{2}(\sigma_i+\sigma_{i+1})}.
\end{equation}
Evaluating for $\sigma_i,\sigma_{i+1}\in\{+1,-1\}$:
\begin{align}
T_{++} &= e^{\beta J+\beta h},\\
T_{+-} = T_{-+} &= e^{-\beta J},\\
T_{--} &= e^{\beta J-\beta h}.
\end{align}
This gives the stated $2\times 2$ matrix.  The partition function is
$Z_N=\Tr T^N$ (standard transfer-matrix result for periodic chains).

\emph{Step 2: Eigenvalues.}
The eigenvalues of $T$ are
\begin{equation}
\lambda_\pm = e^{\beta J}\!\left(\cosh(\beta h)
\pm\sqrt{\sinh^2(\beta h)+e^{-4\beta J}}\right).
\end{equation}
Both are positive for all $\beta,J,h>0$, so $T^N$ is trace class
(trivially, in finite dimension).

\emph{Step 3: LoN representation.}
Apply Proposition~\ref{prop:transferLoN}:
\begin{align}
Z_N &= \Tr T^N = \lambda_+^N+\lambda_-^N\\
&= \int_0^1 e^{-2\pi iu}\,
\det\!\left(1+e^{2\pi iu}T^N\right)du\\
&= \int_0^1 e^{-2\pi iu}\,
\left(1+e^{2\pi iu}\lambda_+^N\right)
\left(1+e^{2\pi iu}\lambda_-^N\right)du.
\end{align}
Expanding the product:
\begin{align}
&= \int_0^1 e^{-2\pi iu}\!\left(1
+e^{2\pi iu}(\lambda_+^N+\lambda_-^N)
+e^{4\pi iu}\lambda_+^N\lambda_-^N\right)du\\
&= 0 + (\lambda_+^N+\lambda_-^N)\int_0^1 du
+ \lambda_+^N\lambda_-^N\int_0^1 e^{2\pi iu}\,du\\
&= \lambda_+^N+\lambda_-^N + 0\\
&= Z_N.\qquad\checkmark
\end{align}
(The $e^{-2\pi iu}$ and $e^{4\pi iu}\cdot e^{-2\pi iu}=e^{2\pi iu}$
terms integrate to zero by orthogonality of characters.)
\end{proof}

\begin{remark}
This provides a concrete answer to the question: ``Can the Ising model
be exactly derived from LoNalogy?''  Yes.  The 1D Ising partition
function is the first holonomy mode of a rank-2 LoN determinant bundle.
The bridge rank $\chi=2$ corresponds to the two spin states $\{\pm 1\}$.
\end{remark}


\section{Boundary-Rank Optimality}\label{sec:optimality}

\begin{theorem}[Boundary-rank optimality]\label{thm:optimality}
For a nearest-neighbour $q$-state model on a strip of width $|B|$
with generic coupling constants, the bridge rank satisfies
\begin{equation}
\boxed{\chi \geq q^{|B|}.}
\end{equation}
In particular, the boundary scaling $\chi=q^{R|B|}$ of
Theorem~\ref{thm:classical} cannot be universally improved.
\end{theorem}

\begin{proof}
\emph{Step 1: Generic transfer matrix.}
For a nearest-neighbour model ($R=1$) on a strip of width $|B|$,
the row-to-row transfer matrix $T_\beta$ is a
$q^{|B|}\times q^{|B|}$ matrix.  For generic coupling constants
(Boltzmann weights), $T_\beta$ has $q^{|B|}$ distinct nonzero
eigenvalues $\lambda_1,\ldots,\lambda_{q^{|B|}}$.

\emph{Step 2: Power-sum argument.}
Suppose there exists an auxiliary operator $A$ of dimension
$\chi<q^{|B|}$ with eigenvalues $\mu_1,\ldots,\mu_\chi$ such that
\begin{equation}
\Tr T_\beta^L = \Tr A^L\qquad\text{for all }L\geq 1.
\end{equation}
This means the power sums agree:
\begin{equation}
p_L(T) := \sum_{j=1}^{q^{|B|}}\lambda_j^L
= \sum_{k=1}^{\chi}\mu_k^L =: p_L(A)
\qquad\text{for all }L\geq 1.
\end{equation}

\emph{Step 3: Newton's identities.}
By Newton's identities, the power sums $\{p_L\}_{L\geq 1}$ uniquely
determine the elementary symmetric polynomials
$\{e_k\}_{k\geq 1}$ of the eigenvalues, and hence determine the
characteristic polynomial.  Specifically, for a set of $n$ values
$\{x_1,\ldots,x_n\}$, the first $n$ power sums $p_1,\ldots,p_n$
determine all elementary symmetric polynomials $e_1,\ldots,e_n$
via the recursive relation:
\begin{equation}
k\,e_k = \sum_{i=1}^{k}(-1)^{i-1}e_{k-i}\,p_i.
\end{equation}

\emph{Step 4: Eigenvalue multiset equality.}
Since $p_L(T)=p_L(A)$ for all $L\geq 1$, and the $T$-eigenvalues
are $q^{|B|}$ distinct nonzero values, the Newton identities applied
to $T$ yield $q^{|B|}$ independent constraints.  For $A$ with
$\chi<q^{|B|}$ eigenvalues, the characteristic polynomial of $A$
has degree $\chi<q^{|B|}$, but must match the power sums of a
degree-$q^{|B|}$ polynomial.  This is impossible unless the
``extra'' eigenvalues of $T$ are zero---but by assumption they are
nonzero.

More precisely: the nonzero eigenvalue multiset of $T_\beta^L$ is
$\{\lambda_j^L\}_{j=1}^{q^{|B|}}$, and that of $A^L$ is
$\{\mu_k^L\}_{k=1}^{\chi}$.  The identity
$\sum\lambda_j^L=\sum\mu_k^L$ for all $L\geq 1$ implies (by the
injectivity of the Laplace/Dirichlet series) that the multisets
$\{\lambda_j\}$ and $\{\mu_k\}$ are equal (where we allow zero
eigenvalues in $A$ to pad).  Since $\{\lambda_j\}$ has $q^{|B|}$
distinct nonzero values, we need $\chi\geq q^{|B|}$.
\end{proof}

\begin{remark}[Optimality]
The transfer-matrix construction (Theorem~\ref{thm:classical}) achieves
$\chi=q^{R|B|}$.  The no-go theorem shows that $\chi\geq q^{|B|}$ is
required for generic nearest-neighbour models.  For $R=1$ (nearest
neighbour), the construction achieves the lower bound exactly:
$\chi=q^{|B|}$.  The transfer class is therefore \emph{optimal} for
generic finite-range local systems.
\end{remark}


%% ============================================================
\part{Quantum Many-Body Extension}
%% ============================================================

\section{Trotter Embedding}\label{sec:trotter}

\begin{theorem}[Quantum systems via Trotter embedding]\label{thm:trotter}
Let $(\mathcal{H},H)$ be a finite-volume quantum many-body system
with $H=H_A+H_B$ (e.g., kinetic + potential, or even/odd sublattice
decomposition).  For any positive integer $m$ (Trotter number), define
the Trotter-approximated partition function
\begin{equation}
Z_m(\beta) := \Tr_{\mathcal{H}}\!\left(
e^{-\beta H_A/m}\,e^{-\beta H_B/m}\right)^m.
\end{equation}
Then:
\begin{enumerate}[label=(\roman*)]
\item $Z_m(\beta)\to Z(\beta)=\Tr e^{-\beta H}$ as $m\to\infty$
      (Trotter product formula).
\item For each finite $m$, $Z_m(\beta)=\Tr T_{\beta,m}^m$ where
      $T_{\beta,m}=e^{-\beta H_A/m}e^{-\beta H_B/m}$ is a transfer
      operator on $\mathcal{H}$.
\item Consequently, for each finite $m$,
\begin{equation}
\boxed{Z_m(\beta) = \int_0^1 e^{-2\pi iu}\,
\det\nolimits_F\!\left(1+e^{2\pi iu}\,T_{\beta,m}^{\,m}\right)du.}
\end{equation}
\end{enumerate}
The quantum partition function is the $m\to\infty$ limit of
LoN determinant bundles.
\end{theorem}

\begin{proof}
\emph{(i)} is the Trotter product formula
($e^{-\beta(H_A+H_B)}=\lim_{m\to\infty}(e^{-\beta H_A/m}e^{-\beta H_B/m})^m$
in trace norm).

\emph{(ii)} Define $T_{\beta,m}:=e^{-\beta H_A/m}e^{-\beta H_B/m}$.
This is a bounded operator on $\mathcal{H}$.  Then
$Z_m(\beta)=\Tr(T_{\beta,m}^m)$ by definition.

\emph{(iii)} Apply Theorem~\ref{thm:universal} (or
Proposition~\ref{prop:transferLoN}) with $A=T_{\beta,m}^m$.
Since $\mathcal{H}$ is finite-dimensional, $T_{\beta,m}^m$ is
trivially trace class.
\end{proof}

\begin{remark}[Quantum = classical + one extra dimension]
The Trotter decomposition maps a $d$-dimensional quantum system to
a $(d+1)$-dimensional classical system with $m$ layers in the
imaginary-time direction.  Each layer is a transfer step.  The
LoN determinant representation at finite $m$ is therefore a
locality-preserving representation on the $(d+1)$-dimensional
lattice, with boundary dimension determined by the transverse
section.

This is the standard quantum-classical correspondence, now given
a LoN determinant interpretation.
\end{remark}

\begin{corollary}[Quantum partition function as limit of LoN bundles]\label{cor:quantumlimit}
For any finite-volume quantum many-body system:
\begin{equation}
\Tr\,e^{-\beta H}
= \lim_{m\to\infty}\int_0^1 e^{-2\pi iu}\,
\det\nolimits_F\!\left(1+e^{2\pi iu}\,T_{\beta,m}^{\,m}\right)du.
\end{equation}
The interchange of limit and integral is justified by the dominated
convergence theorem (since $|\det_F(1+zA)|\leq\det_F(1+|A|)$ for
$|z|=1$, and $\Tr|T_{\beta,m}^m|$ is uniformly bounded in $m$).
\end{corollary}

\begin{proof}
Combine Theorem~\ref{thm:trotter}(i) and (iii).  The dominated
convergence theorem applies because $T_{\beta,m}$ is a product of
bounded positive operators on a finite-dimensional space, and
$\|T_{\beta,m}^m\|_1\leq\|e^{-\beta H_A}\|_1\cdot\|e^{-\beta H_B}\|_1$
uniformly in $m$ (by the Golden--Thompson inequality and trace-norm
submultiplicativity).
\end{proof}


%% ============================================================
\part{Complete Classification}
%% ============================================================

\section{Summary of Results}\label{sec:summary}

\begin{center}
\begin{tabular}{p{5.5cm}lll}
\toprule
Result & Reference & Level\\
\midrule
\multicolumn{3}{l}{\textbf{Gaussian reduction core (Part~I)}}\\
\quad Schur complement = Gaussian elim. & Thm~\ref{thm:KT} & A (exact)\\
\quad Three-sector Schur & Thm~\ref{thm:threesector} & A (exact)\\
\quad Stat.~mech.~as special case & Prop~\ref{prop:statmech} & A (exact)\\
\midrule
\multicolumn{3}{l}{\textbf{Boltzmann, Matsubara, KMS, Gibbs (Parts~II--V)}}\\
\quad Boltzmann from $q$-expansion & Thm~\ref{thm:boltzmann} & A (exact)\\
\quad Matsubara from spin structure & Thm~\ref{thm:matsubara} & A (exact)\\
\quad KMS from Tomita--Takesaki & Thm~\ref{thm:kms} & A (exact)\\
\quad Gibbs state from KMS & Thm~\ref{thm:gibbs} & A (exact)\\
\quad Gibbs variational principle & Thm~\ref{thm:gibbsvar} & A (exact)\\
\quad Hawking temperature & Prop~\ref{prop:hawking} & A (cond.)\\
\midrule
\multicolumn{3}{l}{\textbf{Universal embedding (Part~VI)}}\\
\quad Projected-determinant universality & Thm~\ref{thm:universal} & A (exact)\\
\quad Generating functional & Cor~\ref{cor:generating} & A (exact)\\
\quad Single-fiber insufficiency & Prop~\ref{prop:singlefiber} & A (exact)\\
\midrule
\multicolumn{3}{l}{\textbf{No-go and optimality (Parts~VII--VIII)}}\\
\quad No universal low-rank compression & Thm~\ref{thm:nogo} & A (exact)\\
\quad Classical local $\Rightarrow$ transfer class & Thm~\ref{thm:classical} & A (exact)\\
\quad 1D Ising in LoN determinant class & Cor~\ref{cor:ising} & A (exact)\\
\quad Boundary-rank optimality & Thm~\ref{thm:optimality} & A (exact)\\
\midrule
\multicolumn{3}{l}{\textbf{Quantum extension (Part~IX)}}\\
\quad Trotter embedding & Thm~\ref{thm:trotter} & A (exact)\\
\quad Quantum limit of LoN bundles & Cor~\ref{cor:quantumlimit} & A (exact)\\
\bottomrule
\end{tabular}
\end{center}

\section{Four-Part Classification}\label{sec:fourpart}

The complete verdict on LoNalogy--statistical-mechanics universality is:

\begin{enumerate}[label=\textbf{(\Alph*)}]
\item \textbf{Exact spectral universality: TRUE.}
Every finite-volume Gibbs many-body system is an exact first holonomy
mode of a LoN determinant bundle (Theorem~\ref{thm:universal}).

\item \textbf{Universal low-rank single-fiber compression: FALSE.}
No auxiliary system of dimension $r<D$ can exactly reproduce a generic
$D$-dimensional partition function for all $\beta$
(Theorem~\ref{thm:nogo}).

\item \textbf{Exact local universality: TRUE.}
Finite-range local classical systems are exactly the transfer class,
hence the locality-preserving LoN determinant class
(Theorem~\ref{thm:classical}).

\item \textbf{Boundary-rank optimality: TRUE.}
Generic local models require boundary-rank scaling $\chi\geq q^{|B|}$;
this cannot be universally improved (Theorem~\ref{thm:optimality}).
\end{enumerate}

In one formula:
\begin{equation}
\boxed{\text{LoN determinant universality is true in projected form,
false in universal low-rank fiber form,
and exact/optimal in transfer-boundary form.}}
\end{equation}


%% ============================================================
\section*{Conclusion}
\addcontentsline{toc}{section}{Conclusion}
%% ============================================================

We have established the complete relationship between the LoNalogy
determinant sector and statistical mechanics.  The results form a
logical chain with two branches:

\textbf{Branch 1: Why Gibbs?}
\begin{equation}
X(2)\;\to\; Q=e^{-2\pi\tau_2}\;\to\; e^{-\beta E}
\;\to\;\text{Matsubara}\;\to\;\text{KMS}\;\to\;\text{Gibbs}.
\end{equation}
The Boltzmann factor is the nome of the upper half-plane.  The Gibbs
state is the unique entropy maximiser of the $X(2)$-modular KMS state.
Neither is a postulate.

\textbf{Branch 2: How far does universality go?}
\begin{equation}
\text{Projected bundle: exact}
\;\xrightarrow{\;\text{no-go}\;}
\text{Single fiber: needs full rank}
\;\xrightarrow{\;\text{transfer}\;}
\text{Local systems: optimal.}
\end{equation}
Every Gibbs system embeds into a LoN determinant bundle.  No universal
low-rank compression exists.  Local systems achieve optimal
boundary-rank via the transfer class.  Quantum systems are included
via Trotter embedding.

The four-part classification (A)--(D) is complete: no open problems
remain in the universality question.  The logical chain from $X(2)$
to the full structure of equilibrium statistical mechanics is closed.

In one sentence:

\begin{quote}
\emph{Gibbs is not an axiom.  It is a modular theorem of $X(2)$, and
Boltzmann's exponential is the nome of the upper half-plane.  Every
Gibbs system lives in the projected determinant bundle, and local
systems achieve optimal boundary-rank in the transfer class.}
\end{quote}


%% ============================================================
\section*{Acknowledgements}
\addcontentsline{toc}{section}{Acknowledgements}
%% ============================================================

The author thanks Claude (Anthropic) for assistance with derivation
verification, proof organisation, and manuscript preparation.
The key observation that the Boltzmann factor is the $q$-uniformisation
of $X(2)$ was identified through collaborative analysis of the LoNalogy
determinant sector.  The no-go theorem and transfer-class classification
were developed in the same collaborative session.


%% ============================================================
\begin{thebibliography}{99}

\bibitem{LoNalogyv87}
M.~Nakamura,
\textit{LoNalogy v87.1: Elliptic Modular Framework for Quantum Field Theory
and Millennium Problems},
Zenodo (2026), \href{https://doi.org/10.5281/zenodo.19115338}{doi:10.5281/zenodo.19115338}.

\bibitem{YMmassGap}
M.~Nakamura,
\textit{Yang--Mills Existence and Mass Gap: A Constructive Proof via
Elliptic Modular Geometry and Transfer-Poincar\'e Descent},
Zenodo (2026), \href{https://doi.org/10.5281/zenodo.19183838}{doi:10.5281/zenodo.19183838}.

\bibitem{ParticleCosmo}
M.~Nakamura,
\textit{Standard Model Structure, GUT Scale, and Dark Sector from Elliptic
Modular Geometry: Predictions of LoNalogy},
companion paper (2026).

\bibitem{BratteliRobinson}
O.~Bratteli and D.~W.~Robinson,
\textit{Operator Algebras and Quantum Statistical Mechanics},
Vol.~2, 2nd ed., Springer (1997).

\bibitem{Simon}
B.~Simon,
\textit{Trace Ideals and Their Applications},
2nd ed., Mathematical Surveys and Monographs \textbf{120},
AMS (2005).

\bibitem{Haag}
R.~Haag,
\textit{Local Quantum Physics: Fields, Particles, Algebras},
2nd ed., Springer (1996).

\bibitem{Boltzmann1877}
L.~Boltzmann,
\textit{\"Uber die Beziehung zwischen dem zweiten Hauptsatze der
mechanischen W\"armetheorie und der Wahrscheinlichkeitsrechnung
respektive den S\"atzen \"uber das W\"armegleichgewicht},
Wien.\ Ber.\ \textbf{76}, 373--435 (1877).

\bibitem{Jaynes1957}
E.~T.~Jaynes,
\textit{Information Theory and Statistical Mechanics},
Phys.\ Rev.\ \textbf{106}, 620--630 (1957).

\bibitem{Baxter}
R.~J.~Baxter,
\textit{Exactly Solved Models in Statistical Mechanics},
Academic Press (1982).

\bibitem{Trotter1959}
H.~F.~Trotter,
\textit{On the Product of Semi-Groups of Operators},
Proc.\ Amer.\ Math.\ Soc.\ \textbf{10}, 545--551 (1959).

\end{thebibliography}

\end{document}
